\documentclass[11pt]{letter}
\usepackage[a4paper,margin=1in]{geometry}
\usepackage[T1]{fontenc}
\usepackage{lmodern}
\usepackage{hyperref}
\signature{Author One\\ on behalf of all authors}
\address{Institute/Department\\ National Taiwan University\\ Taipei, Taiwan\\ email@ntu.edu.tw}
\begin{document}
\begin{letter}{The Editors\\ \emph{Briefings in Bioinformatics}\\ Oxford University Press}
\opening{Dear Editors,}

We are pleased to submit our manuscript, ``\emph{How should genomic foundation
models be evaluated for metagenomic taxonomic classification? A controlled
benchmark of pre-training, tokenization, and data scaling},'' for consideration
as a \textbf{Problem-solving Protocol} in \emph{Briefings in Bioinformatics}.

Genomic foundation models (GFMs) are increasingly proposed for metagenomic
tasks, but the community lacks clear, controlled guidance on \emph{which} design
choice actually determines their performance for short-read taxonomic
classification. Published comparisons typically vary model size, data, and
tokenization together, leaving practitioners unable to attribute gains. Our
manuscript addresses this methodological question directly, in the
protocol-and-recommendation style that fits the journal's readership rather than
as a single new-method paper.

We believe the work suits \emph{Briefings in Bioinformatics} for three reasons.
First, its contribution is a \emph{benchmarking protocol and a set of practical
recommendations}, not a single tool---each experiment isolates one factor (data
scale, pre-training, tokenization) so the resulting guidance is causal and
actionable. Second, the central finding is broadly useful and somewhat
counter-intuitive: tokenization, not backbone pre-training or data scale, sets
the ceiling for short-read taxonomy, and a parameter-free 13-mer classifier can
outperform a 498M-parameter pre-trained model. Third, we provide a dual
read-level/sample-level evaluation and a coverage-matched comparison with
Kraken\,2 that yields concrete, use-case-specific recommendations---including the
non-obvious result that a mid-accuracy neural model can exceed Kraken\,2 on
abundance estimation while remaining a weak detector.

We are transparent about the main limitation---the high closed-set accuracy
reflects near-complete reference coverage rather than novel-genome
generalization---and we frame it as a direction for the field rather than a
solved problem; an out-of-genome generalization experiment is in progress and
can be added during revision if the editors consider it essential.

This manuscript is original, is not under consideration elsewhere, and all
authors approve its submission. We disclose in the manuscript that an AI
assistant was used to help draft and edit the text and organize results; all
scientific claims and reported numbers were verified by the authors. We declare
no competing interests. We would be glad to suggest suitable reviewers with
expertise in genomic language models and metagenomic classification upon
request.

\closing{Sincerely,}
\end{letter}
\end{document}
